% Supplementary-material fragment: verifier training, rendering, and reproducibility.
% This file is intended to be included with \input{...}; it has no preamble.
% TODO--AUTHOR items must be resolved against the frozen release artifacts.
\providecommand{\todovalue}{\textbf{[TODO--AUTHOR]}}
\providecommand{\todoauthor}{\textbf{TODO--AUTHOR:}}

\section{Verifier Training, Rendering, and Reproducibility Details}
\label{sec:supp-verifier}

\paragraph{Scope.}
This section documents the parameter-efficient verifier runs based on
low-rank adaptation (LoRA), together with the exact prompt construction,
training objective, checkpoint rule, and deterministic label-token
inference path.  The settings below define a reference LoRA recipe and its
dataset-specific variants.  They should not be interpreted as a claim that
one numerically identical configuration applies to every experiment;
dataset-specific differences are stated explicitly.

\subsection{From an Ordered Trace to the Verifier Input}
\label{sec:supp-verifier-renderer}

\paragraph{Prompt-visible evidence prefix.}
The selector first produces an ordered trace
$\mathcal{T}=(u_1,\ldots,u_T)$.  The renderer consumes an
order-preserving prefix rather than reranking the trace.  Under the
\texttt{minmax} policy, it appends evidence in trace order and consults the
stored \texttt{target\_resolved} flag after each appended step.  This flag is
the implementation record that the step's replayed resolved-atom rate has
reached its configured target.  Writing $q_t\in\{0,1\}$ for the flag and
$H=\min(T,K_{\max})$, let
\[
\mathcal Q=\{t:K_{\min}\leq t\leq H,\ q_t=1\}.
\]
The visible length is
\begin{equation}
K^\star =
\begin{cases}
\min\mathcal Q,&\mathcal Q\neq\varnothing,\\
H,&\mathcal Q=\varnothing.
\end{cases}
\label{eq:supp-visible-prefix}
\end{equation}
Candidate exhaustion can therefore yield fewer than $K_{\min}$ units.
The reference count policies are minmax$(5,10)$ for LIAR-RAW and RAWFC
and minmax$(9,9)$ for SciFact.  Equation~\ref{eq:supp-visible-prefix}
describes the verifier-visible prefix policy; it is distinct from the
selector's construction of the longer audit ordering.

\paragraph{Trace rendering.}
For the \texttt{mrec\_min} style, each selected candidate is aligned to
its trace step.  The renderer obtains the step-specific cue, falling back
to an atom-derived or generic verification cue when the step does not
provide one, and serializes the unit exactly as
\begin{lstlisting}
Check: <STEP-SPECIFIC CUE>
<ORIGINAL EVIDENCE TEXT>
\end{lstlisting}
The outer evidence formatter then enumerates these strings as
\texttt{[1]}, \texttt{[2]}, and so on, preserving trace order.  A cue is
therefore attached to each visible evidence unit; the same atom cue may
appear for more than one unit.  Trace metadata such as operation,
relation, directness, covered-atom identifiers, and before/after atom
states remains in the audit artifact but is not serialized in the
\texttt{mrec\_min} verifier prompt.

\paragraph{Complete prompt template.}
The system message is instantiated with the dataset task name:
\begin{lstlisting}
You are a careful fact-checking assistant for <TASK-NAME> claims. Classify claims using only the claim and retrieved evidence supplied by the user.
\end{lstlisting}
The user message is:
\begin{lstlisting}
Classify the claim into exactly one <TASK-NAME> label.

Labels:
<SCHEMA-SPECIFIC LABEL DEFINITIONS>

Rules:
- Use the retrieved evidence as the primary source.
- Do not invent facts not supported by the evidence.
- Respond with exactly one line: Label: <a single letter from A-<LAST>>

Claim:
<CLAIM>

Evidence:
[1] Check: <CUE-1>
<EVIDENCE-1>
[2] Check: <CUE-2>
<EVIDENCE-2>
...
[K] Check: <CUE-K>
<EVIDENCE-K>
\end{lstlisting}
If no evidence is available, the evidence block is the literal string
\texttt{(no evidence available)}.  The placeholder
\texttt{<SCHEMA-SPECIFIC LABEL DEFINITIONS>} expands as follows.

\begin{table*}[t]
\centering
\small
\setlength{\tabcolsep}{4pt}
\renewcommand{\arraystretch}{1.08}
\caption{Exact label-token schemas and label descriptions inserted into the verifier prompt.}
\label{tab:supp-label-schemas}
\begin{tabular}{p{0.13\textwidth}p{0.10\textwidth}p{0.69\textwidth}}
\hline
\textbf{Dataset} & \textbf{Token} & \textbf{Label and prompt definition}\\
\hline
LIAR-RAW & A & pants-fire: completely false and implausible\\
         & B & false: false based on the available evidence\\
         & C & barely-true: mostly false, with only a small element of truth\\
         & D & half-true: partly true and partly false\\
         & E & mostly-true: mostly true, with minor missing context or caveats\\
         & F & true: accurate based on the available evidence\\
\hline
RAWFC & A & false: false based on the available evidence\\
      & B & half: partly true and partly false\\
      & C & true: accurate based on the available evidence\\
\hline
SciFact & A & support: the scientific claim is supported by the retrieved abstract evidence\\
        & B & contradict: the scientific claim is contradicted by the retrieved abstract evidence\\
        & C & nei: the retrieved abstract evidence is insufficient to support or contradict the scientific claim\\
\hline
\end{tabular}
\end{table*}

The system and user messages are passed to the base tokenizer's native
chat template with \texttt{add\_generation\_prompt=true}; no handwritten
common chat wrapper is substituted across backbones.  For the
MistralCommon tokenizer path, the tokenized chat-template output is stored
with the built example and reused directly.  The training dataset then
appends the tokenized prefix \texttt{Label:} after the assistant-generation
header.  Thus, the supervised decision is the single token immediately
following \texttt{Label:}, rather than a generated explanation or a
multi-token answer string.

\subsection{Restricted Label-Token Training and Inference}
\label{sec:supp-label-token}

Let $\mathcal{Y}=(y_0,\ldots,y_{C-1})$ denote the ordered label schema and
let $g(y_j)$ be its letter token.  Because \texttt{Label:} does not end in
whitespace, the implementation tokenizes the choice text as
\texttt{\char32 A}, \texttt{\char32 B}, and so forth.  It verifies at run
startup that every choice text maps to exactly one tokenizer token.
Given the final hidden position after the prompt and \texttt{Label:}, only
the logits of these $C$ permitted tokens are retained:
\begin{equation}
p_{nj} =
\frac{\exp z_{n,g(y_j)}}{\sum_{k=0}^{C-1}\exp z_{n,g(y_k)}}.
\label{eq:supp-restricted-softmax}
\end{equation}

\paragraph{Class-weighted cross entropy.}
For class weights $w_j>0$, the implemented batch loss normalizes by the
sum of the sample weights:
\begin{equation}
\mathcal{L}_{\mathrm{wCE}}
=
\frac{
\sum_{n=1}^{N} w_{y_n}\bigl[-\log p_{n,y_n}\bigr]
}{
\sum_{n=1}^{N} w_{y_n}
}.
\label{eq:supp-weighted-ce}
\end{equation}
The LIAR-RAW LoRA configuration uses weights
$(1.2,1.0,1.5,1.0,1.0,1.8)$ in the ordered class list
(pants-fire, false, barely-true, half-true, mostly-true, true).
The reference RAWFC and SciFact LoRA configurations use unit class weights.

\paragraph{Optional ordinal auxiliary loss.}
The implementation optionally exploits the natural LIAR-RAW label order.
With $r(y_j)=j$, define the normalized ordinal distance
\[
d_{nj}=\frac{|r(y_j)-r(y_n)|}{C-1}.
\]
The auxiliary loss is
\begin{equation}
\mathcal{L}_{\mathrm{ord}}
=
\frac{
\sum_{n=1}^{N}w_{y_n}\sum_{j=0}^{C-1}p_{nj}d_{nj}
}{
\sum_{n=1}^{N}w_{y_n}
},
\label{eq:supp-ordinal}
\end{equation}
and optimizes
\begin{equation}
\mathcal{L}_{\mathrm{ver}}
=
\mathcal{L}_{\mathrm{wCE}}
+\alpha_t\mathcal{L}_{\mathrm{ord}}.
\label{eq:supp-verifier-loss}
\end{equation}
For the reference LIAR-RAW and RAWFC LoRA configurations, $\alpha=0.2$ and
$\alpha_t=\alpha\min(1,t/(0.3T_{\mathrm{train}}))$, i.e., the auxiliary
coefficient is linearly warmed up over the first 30\% of planned update
steps.  The SciFact configuration disables this term.  In every case, the
resolved run configuration controls whether the option is active.

\paragraph{Deterministic decoding.}
At inference time, the prediction is
\begin{equation}
\widehat{y}_n
=
\arg\max_{y_j\in\mathcal{Y}}z_{n,g(y_j)}.
\label{eq:supp-label-argmax}
\end{equation}
No autoregressive generation, beam search, text parser, sampling
temperature, or ``thinking'' mode participates in this decision.  The
human-readable string \texttt{Label: X} in prediction artifacts is
assembled after the argmax.  Unless a run is explicitly marked otherwise,
label-prior logit adjustment is disabled.

\subsection{Reference LoRA Optimization Recipe}
\label{sec:supp-lora-recipe}

The base model parameters are frozen and the optimizer updates only LoRA
parameters~\citep{Hu2022LoRA}.  LoRA is applied to
\texttt{q\_proj}, \texttt{k\_proj}, \texttt{v\_proj},
\texttt{o\_proj}, \texttt{gate\_proj}, \texttt{up\_proj}, and
\texttt{down\_proj}, with rank $r=16$, scaling $\alpha=32$, dropout
$0.1$, no trainable bias, and no additional modules-to-save.  This gives
41,943,040 trainable parameters out of 8,072,204,288 (0.5196\%) for the
archived Llama-3.1-8B setup and 53,608,448 out of 8,971,634,688
(0.5975\%) for the archived Ministral-3-8B setup.

The common recipe uses BF16, AdamW, weight decay $0.01$, a 3\% linear
warm-up, a cosine-with-restarts schedule with two cycles, gradient clipping
at $1.0$, gradient checkpointing, longest-in-batch dynamic padding, length
bucketed training, and DeepSpeed ZeRO-2.  The maximum input length is
1,024 tokens.  With four data-parallel processes, per-device batch size
one and four accumulation steps give an effective batch size of 16.
The epoch cap is 12 and the canonical trainer seed is 42.  The archived
LIAR-RAW logs requested FlashAttention~2 but recorded that it was
unavailable and therefore used the model's default attention
implementation.

\begin{table}[t]
\centering
\scriptsize
\setlength{\tabcolsep}{2pt}
\renewcommand{\arraystretch}{1.08}
\caption{Dataset-specific settings for the reference LoRA verifier recipe. No performance values are reported in this table.}
\label{tab:supp-lora-settings}
\begin{tabular}{lrrrr}
\hline
\textbf{Dataset} & \textbf{LR} & \textbf{Eval/save} & \textbf{Patience} & \textbf{Prefix}\\
\hline
LIAR-RAW & $2{\times}10^{-5}$ & 100 & 8 / 12$^\dagger$ & minmax$(5,10)$\\
RAWFC & $1{\times}10^{-5}$ & 50 & 8 & minmax$(5,10)$\\
SciFact & $2{\times}10^{-5}$ & 100 & 12 & minmax$(9,9)$\\
\hline
\end{tabular}
\vspace{2pt}

\footnotesize $^\dagger$Llama / Ministral in the archived LIAR-RAW
reference configurations.  Intervals and patience are measured in
optimizer-update evaluations, not epochs.
\end{table}

\paragraph{Checkpoint selection.}
For an ordinary single-arm run, each evaluation checkpoint is ranked by
validation macro-F1.  When an experiment provides a valid matched-arm
evaluation, the implementation instead uses arm-balanced macro-F1.
Scores equal to within an absolute tolerance of $10^{-12}$ are broken
first by lower validation mean cross entropy and then by the earlier
optimizer step.  The best adapter and the final adapter are saved
separately.  Early-stopping patience counts consecutive evaluation events
without an improvement under this same ordering; it is not a count of
epochs.

\subsection{Length Budgeting and Truncation}
\label{sec:supp-truncation}

For the reference dataset runs in
Table~\ref{tab:supp-lora-settings}, prompt construction and trainer
tokenization enforce the same 1,024-token limit.  Before rendering, the
prefix policy chooses the evidence units
described in Equation~\ref{eq:supp-visible-prefix}.  The prompt builder
then reserves the tokenizer-specific target allowance, including the
end-of-sequence allowance used by the builder, and defines
\begin{equation}
B_{\mathrm{prompt}}=L_{\max}-L_{\mathrm{target\ reserve}}.
\end{equation}
If the rendered prompt exceeds this budget, complete evidence units are
removed from the \emph{tail}, preserving the relative order of all retained
units.  The prompt is rebuilt and re-tokenized after every removal.  If a
single remaining unit still does not fit, the generic fallback performs a
binary search for the longest token prefix of that evidence text that
fits; if even that is impossible, it emits the no-evidence prompt.

Built examples set \texttt{preserve\_prompt\_prefix=true}.  Consequently,
the label-token dataset does not silently left-truncate an overlength
protected prompt after appending \texttt{Label:}; it raises an error and
requires the build-time evidence budget to be corrected.  The frozen
release should report the following diagnostics for each split:
\begin{itemize}
    \item number and fraction of rows with tail-unit removal;
    \item evidence count before and after removal;
    \item number of rows invoking within-unit text truncation;
    \item number of rows ending below $K_{\min}$ because of exhaustion or
    length constraints; and
    \item prompt-token percentiles under the exact verifier tokenizer.
\end{itemize}
\noindent\todoauthor populate these diagnostics from the final
archived build manifests and add their SHA-256 checksums.

\subsection{Reproducibility Ledger}
\label{sec:supp-repro-ledger}

\paragraph{Immutable model and artifact snapshots.}
Public model names alone do not identify immutable weights or chat
templates.  The release manifest should freeze the following fields.

\begin{table*}[t]
\centering
\small
\setlength{\tabcolsep}{3pt}
\renewcommand{\arraystretch}{1.08}
\caption{Model, tokenizer, and adapter snapshot ledger.}
\label{tab:supp-model-snapshots}
\begin{tabular}{p{0.18\textwidth}p{0.28\textwidth}p{0.16\textwidth}p{0.28\textwidth}}
\hline
\textbf{Role} & \textbf{Public identifier} & \textbf{Immutable revision} & \textbf{Checksums / release record}\\
\hline
Verifier backbone &
\path{mistralai/Ministral-3-8B-Instruct-2512} &
\todovalue &
Weight index or consolidated-weight SHA-256; tokenizer and
\texttt{chat\_template.jinja} SHA-256: \todovalue\\
Verifier backbone &
\path{meta-llama/Meta-Llama-3.1-8B-Instruct} &
\todovalue &
Weight-index SHA-256; tokenizer and tokenizer-config SHA-256:
\todovalue\\
Cross-verifier backbone &
\path{Qwen/Qwen3-4B-Instruct-2507} &
\todovalue &
Weight-index SHA-256; tokenizer, tokenizer-config, and chat-template
SHA-256: \todovalue\\
LoRA adapters &
One adapter directory per frozen run &
Best optimizer step: \todovalue &
\texttt{adapter\_config.json}, adapter-weight, resolved-config, and
prediction-file SHA-256: \todovalue\\
Prompt data &
Built train/validation/test JSONL files &
Schema version: \todovalue &
Row counts, prompt-template hash, and JSONL SHA-256:
\todovalue\\
\hline
\end{tabular}
\end{table*}

\paragraph{Software and execution environment.}

\begin{table*}[t]
\centering
\small
\setlength{\tabcolsep}{4pt}
\renewcommand{\arraystretch}{1.08}
\caption{Software and hardware information to freeze with the release. ``Recorded constraint'' is not a substitute for the exact installed version.}
\label{tab:supp-software-hardware}
\begin{tabular}{p{0.22\textwidth}p{0.30\textwidth}p{0.36\textwidth}}
\hline
\textbf{Item} & \textbf{Recorded constraint / behavior} & \textbf{Exact released value}\\
\hline
Code & Python implementation in this repository &
Git commit, dirty-state flag, and patch archive:
\todovalue\\
Python / OS & Python 3.10+; Linux &
Python, distribution, and kernel versions:
\todovalue\\
PyTorch stack & PyTorch 2.6.0 with CUDA 12.4 wheels is pinned in
\texttt{requirements.txt} &
Installed PyTorch, CUDA runtime, cuDNN, and NCCL versions:
\todovalue\\
Training libraries &
Transformers $\geq4.39$, PEFT $\geq0.12$, Accelerate $\geq1.8$,
DeepSpeed $\geq0.15$ &
Exact versions and environment lock file:
\todovalue\\
Attention backend &
Default model attention in the audited LIAR-RAW LoRA logs after the
FlashAttention~2 availability check failed &
Confirm backend independently for every released run:
\todovalue\\
Numeric modes & BF16 training; TF32 enabled on supported CUDA devices &
Determinism flags and any environment variables:
\todovalue\\
Accelerators & NVIDIA L20 family reported by the experiment notes &
Number of hosts, GPUs per host, GPU memory, driver, and topology:
\todovalue\\
Host resources & Not frozen in current artifacts &
CPU model/count, RAM, local-storage type, and interconnect:
\todovalue\\
\hline
\end{tabular}
\end{table*}

\paragraph{Compute and monetary cost.}
Wall time must be measured from the frozen run logs rather than inferred
from the epoch cap, because early stopping and resumed runs change the
actual cost.  GPU-hours are computed as the number of simultaneously
allocated GPUs multiplied by elapsed wall-clock hours.

\begin{table*}[t]
\centering
\small
\setlength{\tabcolsep}{3pt}
\renewcommand{\arraystretch}{1.08}
\caption{Compute ledger for the released LoRA verifier runs. This table deliberately contains no task-performance values.}
\label{tab:supp-compute-cost}
\begin{tabular}{p{0.17\textwidth}p{0.17\textwidth}p{0.12\textwidth}p{0.12\textwidth}p{0.12\textwidth}p{0.18\textwidth}}
\hline
\textbf{Dataset} & \textbf{Backbone} & \textbf{GPU count/type} &
\textbf{Wall time} & \textbf{GPU-hours} & \textbf{Peak memory / cost}\\
\hline
LIAR-RAW & Ministral-3-8B & \todovalue &
\todovalue & \todovalue &
\todovalue\\
LIAR-RAW & Llama-3.1-8B & \todovalue &
\todovalue & \todovalue &
\todovalue\\
RAWFC & Ministral-3-8B & \todovalue &
\todovalue & \todovalue &
\todovalue\\
RAWFC & Llama-3.1-8B & \todovalue &
\todovalue & \todovalue &
\todovalue\\
SciFact & Ministral-3-8B & \todovalue &
\todovalue & \todovalue &
\todovalue\\
LIAR-RAW cross-verifier & Qwen3-4B & \todovalue &
\todovalue & \todovalue &
\todovalue\\
LIAR-RAW cross-verifier & Llama-3.1-8B & \todovalue &
\todovalue & \todovalue &
\todovalue\\
\hline
\end{tabular}
\end{table*}

\noindent\todoauthor state the accounting basis for monetary
cost (institutional hardware, cloud list price, or amortized internal
rate), report verifier-training and verifier-inference costs separately,
and indicate whether failed, calibration, and ablation runs are included.
If energy or carbon estimates are reported, also record the measurement
tool, power-usage effectiveness assumption, grid region, and date.

\FloatBarrier
